\documentclass[12pt]{article}
\usepackage{amsmath,amssymb}

\begin{document}
\section*{{2024 AIME I Problems/Problem 4}}
Source: AoPS Wiki

\subsection*{Solution}
This is a conditional probability problem. Bayes' Theorem states that \[P(A|B)=\dfrac{P(B|A)\cdot P(A)}{P(B)}\] in other words, the probability of $A$ given $B$ is equal to the probability of $B$ given $A$ times the probability of $A$ divided by the probability of $B$ . In our case, $A$ represents the probability of winning the grand prize, and $B$ represents the probability of winning a prize. Clearly, $P(B|A)=1$ , since by winning the grand prize you automatically win a prize. Thus, we want to find $\dfrac{P(A)}{P(B)}$ . Let us calculate the probability of winning a prize. We do this through casework: how many of Jen's drawn numbers match the lottery's drawn numbers? To win a prize, Jen must draw at least $2$ numbers identical to the lottery. Thus, our cases are drawing $2$ , $3$ , or $4$ numbers identical. Let us first calculate the number of ways to draw exactly $2$ identical numbers to the lottery. Let Jen choose the numbers $a$ , $b$ , $c$ , and $d$ ; we have $\dbinom42$ ways to choose which $2$ of these $4$ numbers are identical to the lottery. We have now determined $2$ of the $4$ numbers drawn in the lottery; since the other $2$ numbers Jen chose can not be chosen by the lottery, the lottery now has $10-2-2=6$ numbers to choose the last $2$ numbers from. Thus, this case is $\dbinom62$ , so this case yields $\dbinom42\dbinom62=6\cdot15=90$ possibilities. Next, let us calculate the number of ways to draw exactly $3$ identical numbers to the lottery. Again, let Jen choose $a$ , $b$ , $c$ , and $d$ . This time, we have $\dbinom43$ ways to choose the identical numbers and again $6$ numbers left for the lottery to choose from; however, since $3$ of the lottery's numbers have already been determined, the lottery only needs to choose $1$ more number, so this is $\dbinom61$ . This case yields $\dbinom43\dbinom61=4\cdot6=24$ . Finally, let us calculate the number of ways to all $4$ numbers matching. There is actually just one way for this to happen, as $\dbinom44=1$ . In total, we have $90+24+1=115$ ways to win a prize. The lottery has $\dbinom{10}4=210$ possible combinations to draw, so the probability of winning a prize is $\dfrac{115}{210}$ . There is actually no need to simplify it or even evaluate $\dbinom{10}4$ or actually even know that it has to be $\dbinom{10}4$ ; it suffices to call it $a$ or some other variable, as it will cancel out later. However, let us just go through with this. The probability of winning a prize is $\dfrac{115}{210}$ . Note that the probability of winning a grand prize is just matching all $4$ numbers, which we already calculated to have $1$ possibility and thus have probability $\dfrac1{210}$ . Thus, our answer is $\dfrac{\frac1{210}}{\frac{115}{210}}=\dfrac1{115}$ . Therefore, our answer is $1+115=\boxed{116}$ . ~Technodoggo

\end{document}
